\section{Related Work}
\paragraph{Retrieval-augmented generation and adaptive retrieval.}
Retrieval-augmented generation combines parametric generation with non-parametric evidence access \citep{lewis2020rag}. Dense passage retrieval established a strong neural retrieval paradigm for open-domain question answering \citep{karpukhin2020dpr}, while later systems interleave retrieval with reasoning, reflection, or critique \citep{trivedi2023interleaving,asai2024selfrag}. These methods motivate interactive evidence acquisition, but many assume a document collection rather than a text-rich graph whose structure constrains valid answers.

\paragraph{Language-model agents and tool use.}
Tool-using agents such as ReAct combine reasoning traces with environment actions \citep{yao2023react}. Reflexion, Self-Refine, and Toolformer explore variants of feedback, iterative improvement, and learned tool use \citep{shinn2023reflexion,madaan2023selfrefine,schick2023toolformer}. In principle, such methods suggest that a retrieval controller should improve after seeing observations. Our results identify a boundary condition where observation feedback can make the process more measurable while still degrading final retrieval when the next-action policy is not calibrated.

\paragraph{Graph retrieval over textual and relational knowledge bases.}
Knowledge-graph reasoning methods such as Think-on-Graph expose graph traversal to language models \citep{sun2023thinkongraph}. \stark is the most relevant benchmark for our setting because it evaluates retrieval over semi-structured textual and relational knowledge bases across product search, academic paper search, and precision medicine \citep{wu2024stark}. Autofocus or AF-Retriever studies multi-hop retrieval over semi-structured knowledge bases by combining entity search, LLM-generated Cypher queries, node-set joins, vector search, and reranking \citep{boer2025afretriever}. GraphRAFT explores retrieval-augmented fine-tuning for graph databases \citep{clemedtson2025graphraft}. ARK is the closest architectural reference and it gives a language model two graph tools, global lexical search with one-hop neighborhood exploration, so that the model can balance breadth-oriented discovery with depth-oriented traversal \citep{polonuer2026ark}. We build on the same broad premise but focus on trace-conditioned failure analysis rather than claiming a new state-of-the-art retriever.

\paragraph{Personalization before retrieval.}
Personalized retrieval work argues that user-specific semantics should influence retrieval before documents are fetched. Personalize Before Retrieve (PBR) performs personalized query expansion using style-aware pseudo feedback and corpus-structure alignment \citep{zhang2025pbr}. PersonaBench evaluates whether systems can answer questions requiring access to synthetic personal information \citep{tan2025personabench}. Our experiments do not introduce real user profiles; the profile in \tb is query-derived. The relevance to personalization is methodological as a user profile is another pre-retrieval signal, and our results show that adding such a signal globally to an agent prompt is not a substitute for learning when it should change the retrieval action.

\paragraph{Failure diagnostics for agents.}
Final answer metrics can hide process failures as an agent may make redundant calls, chase empty branches, or recover by chance after an inefficient trajectory. Agent evaluation work therefore increasingly emphasizes trajectories, reproducible failure triggers, and process-level measurements \citep{liu2023agentbench,xie2024osworld}. \tb follows this instrumentation-first view. We pair standard retrieval metrics with trajectory counters so that failure is attributed to observable behavior rather than only to aggregate score changes.
